% !TEX root = ../main.tex

\section{Silnik Stirlinga - sprawność w punkcie maksymalnej mocy}
W przeciwieństwie do wcześniejszych cykli w tym ciepło jest wymieniane podczas każdej przemiany. Przemiany izotermiczne zachodzą odpowiednio w temperaturach \(T_{hw}\) oraz \(T_{cw}\), podczas gdy zbiorniki cieplne znajdują się w temperaturach \(T_{h}\) oraz \(T_{c}\). Różnice te oznaczamy podobnie jak w przypadku silnika Carnota \cite{endo_stirling}:
\begin{align}
	T_h - T_{hw} &= x > 0, & T_{cw} - T_c &= y > 0.
\end{align}
Ciepła w przemianach izochorycznych są sobie równe co do wielkości i przeciwne co do znaku.

\begin{figure}[htp]
	\centering
	\tikzsetnextfilename{pv_silnik_endo_stirling_70}
\begin{tikzpicture}
	\begin{axis}[width=0.7\textwidth,
		xlabel=$V$,
		ylabel=$p$,
		axis x line=bottom,
		axis y line=middle,
		xmin=0,
		xmax=1.2,
		ymin=0,
		ymax=1.2,
		ticks=none,
		every axis x label/.style={at={(current axis.right of origin)},anchor=north east},
		every axis y label/.style={at={(current axis.above origin)},anchor=north east},
		declare function={
			r = 2;
			tau = 2;
			V1 = 1/r;
			V2 = 1;
			p1 = 1/r * 1/tau;
			p2 = 1/tau;
			p3 = 1;
			p4 = 1/r;
		}]
		\addplot[mark=*] coordinates {(V2,p1)} node[anchor=north west]{$1$};
		\addplot[mark=*] coordinates {(V1,p2)} node[anchor=north east]{$2$};
		\addplot[mark=*] coordinates {(V1,p3)} node[anchor=south east]{$3$};
		\addplot[mark=*] coordinates {(V2,p4)} node[anchor=south west]{$4$};
		\addplot[domain=V1:V2] {1/r * 1/tau / x} node [pos=0.4, below, sloped] {$T_{cw}$};
		\addplot[mark=none] coordinates {(V1, p2) (V1, p3)};
		\addplot[domain=V1:V2] {1/r / x}         node [pos=0.5, above, sloped] {$T_{hw}$};
		\addplot[mark=none] coordinates {(V2, p4) (V2, p1)};
		\addplot[mark=none, dashed] coordinates {(V2, p1) (V2, 0)};
		\addplot[mark=none, dashed] coordinates {(V1, p2) (V1, 0)};
		\draw[line width=0.3mm, -{Stealth[length=5mm, open]}] (0.4,0.7) node[anchor=east]{$Q_{2-3}$} -- (0.5,0.7);
		\draw[line width=0.3mm, -{Stealth[length=5mm, open]}] (1,0.35) -- (1.1,0.35) node[anchor=north]{$Q_{4-1}$};
		\draw[line width=0.3mm, -{Stealth[length=5mm, open]}] (0.85,0.7) node[anchor=south west]{$Q_{3-4}$} -- (0.8,0.6);
		\draw[line width=0.3mm, -{Stealth[length=5mm, open]}] (0.8,0.3) -- (0.75,0.2) node[anchor=north west]{$Q_{1-2}$};
		
		\addplot[domain=V1:V2, dashed] {0.7*1/r * 1/tau / x} node [pos=0.4, below, sloped] {$T_c$};
		\addplot[domain=1.1*V1:V2, dashed] {1.2*1/r / x}         node [pos=0.5, above, sloped] {$T_h$};
	\end{axis}
\end{tikzpicture}

	\caption{Wykres cyklu Stirling w warunkach endoodwracalnych}
\end{figure}

Rozpatrujemy cykl z gazem doskonałym jako czynnikiem roboczym, dla którego energia wewnętrzna jest jedynie funkcją temperatury \cite{b:kondepudi_prigogine_modern_thermodynamics}. Podczas przemian izotermicznych ciepło i praca spełniają:
\begin{align}
	Q_{1-2} &= - W_{1-2} = \int_{V_1}^{V_2} p\d{V} = nRT_{cw}\ln\pab{\frac{V_2}{V_1}}, \\
	Q_{3-4} &= - W_{3-4} = \int_{V_2}^{V_1} p\d{V} = nRT_{hw}\ln\pab{\frac{V_1}{V_2}}.
\end{align}
Zmiany entropii dla nich zgodnie z równaniem \eqref{eq:def_entropia} spełniają:
\begin{equation}\label{eq:endo_stirling_deltaS}
	\Delta S_{1-2} + \Delta S_{3-4} = \frac{Q_{1-2}}{T_{cw}} + \frac{Q_{3-4}}{T_{hw}} = 0.
\end{equation}
Po pełnym cyklu wracamy do stanu początkowego. Entropia czynnika roboczego jest funkcją stanu, więc w pełnym cyklu jej zmiana jest zerowa:
\begin{equation}
	\Delta S = \Delta S_{1-2} + \Delta S_{2-3} + \Delta S_{3-4} + \Delta S_{4-1} = 0.
\end{equation}
Uwzględniając \eqref{eq:endo_stirling_deltaS} otrzymujemy:
\begin{equation}
	\Delta S_{2-3} + \Delta S_{4-1} = 0.
\end{equation}
Oznacza to, że dwie przemiany izochoryczne w analizie są identyczne jak dwie przemiany izentropowe cyklu Carnota. Ciepło tych przemian nie jest też brane pod uwagę, ponieważ wymieniane jest ono z wewnętrznym regeneratorem, a nie z otoczeniem. Problem ten sprowadza się więc do analizy cyklu Carnota w punkcie maksymalnej mocy i nie będzie powtórzona. Wyniki końcowe są następujace \cite{endo_stirling}:
\begin{gather}
	\eta = 1 - \sqrt{\frac{T_c}{T_h}}, \\
	P = \frac{\alpha_h \alpha_c}{\zeta} \pab{\frac{\sqrt{T_h} - \sqrt{T_c}}{\sqrt{\alpha_h} + \sqrt{\alpha_c}}}^2.
\end{gather}

Wynik ten jest oparty na idealnym działaniu regeneratora. Wymaga to, żeby regenerator zmieniał swoją temperaturę tak, aby była równa temperaturze czynnika roboczego.
